\chapter[Band 41: Abelsche Gruppen]{Band 41 auf einen Blick}
\noindent\textbf{Abelsche Gruppen}\par\medskip
Dieser Band ergänzt die Gruppenaxiome aus Band~40 um Kommutativität.
Die Differenzenkonstruktion setzt ein kommutatives kürzbares Monoid voraus.

\begin{BandUebersicht}
\textbf{Abelsche Gruppe}

\UebFormel{\begin{aligned}&\AbelianGroupAlg G\star e\ \leftrightarrow\\&\GroupAlg G\star e\land\\&\forall x,y\in G\,(x\star y=y\star x).\end{aligned}}
\UebQuellen{\FormulaRefAuto{AbelianGroupDef}[def]}
&
\textbf{Morphismen}
Quelle und Ziel sind abelsche Gruppen. Ein Homomorphismus erhält die Operation; ein Isomorphismus ist zusätzlich bijektiv.
\UebFormel{\AbelianGroupIso f{G,\star,e}{H,\diamond,u}.}
\UebQuellen{\FormulaRefAuto{AbelianGroupHomDef}[def];
\FormulaRefAuto{AbelianGroupIsoDef}[def]} \\
\addlinespace[.8ex]

\textbf{Formale Differenzen}

\UebFormel{\begin{aligned}&(a,b)\sim(c,d)\ \leftrightarrow a\star d=b\star c,\\&K(M)=(M\times M)/{\sim},\\&[a,b]\oplus[c,d]=[a\star c,b\star d].\end{aligned}}
\UebQuellen{\FormulaRefAuto{GroupCompletionPairRelationDef}[def];
\FormulaRefAuto{GroupCompletionCarrierDef}[def];
\FormulaRefAuto{GroupCompletionOperationDef}[def]}
&
\textbf{Gruppenvervollständigung}

\UebFormel{\begin{aligned}&[a,b]=[c,d]\ \leftrightarrow a\star d=b\star c,\\&\AbelianGroupAlg{K(M)}\oplus{[e,e]}.\end{aligned}}
\UebQuellen{\FormulaRefAuto{GroupCompletionClassEquality}[thm];
\FormulaRefAuto{GroupCompletionClassInverse}[thm];
\FormulaRefAuto{GroupCompletionAbelianGroup}[thm]} \\
\addlinespace[.8ex]

\textbf{Kanonische Einbettung}

\UebFormel{\iota(a)=[a,e].}
\UebQuellen{\FormulaRefAuto{GroupCompletionEmbeddingDef}[def]}
&
\textbf{Einbettungssatz}

\UebFormel{\begin{aligned}&\iota:M\inj K(M),\quad\iota(e)=[e,e],\\&\iota(a\star b)=\iota(a)\oplus\iota(b).\end{aligned}}
\UebQuellen{\FormulaRefAuto{GroupCompletionEmbeddingTheorem}[thm]} \\
\addlinespace[.8ex]

\textbf{Ganzzahlige Addition}
Die Rechengesetze stammen aus der Quotientenkonstruktion in Band~17.
\UebFormel{\IntClass a b+\IntClass c d=\IntClass{a+c}{b+d}.}
\UebQuellen{\FormulaRefAuto{IntegerAdditionDef}[def]}
&
\textbf{Additives Grundbeispiel}

\UebFormel{\AbelianGroupAlg{\mathbb Z}{+}{\IntZero}.}
\UebQuellen{\FormulaRefAuto{IntegerAdditiveAbelianGroup}[thm]} \\
\addlinespace[.8ex]
\end{BandUebersicht}
